\scanpage{15}
Supposons $X$ une \uline{surface}.

Notons maintenant que $J'$ a une polarisation
canonique, héritée de celle de $J$. \deleted{\unsure{Soit|Car} $J'_0$
une sous [En effet, \unsure{la|le} \unsure{du|dual de} $J$ étant sa
$J\to$ propre duale, $\unsure{\check J'|J'|\hat J'}$ la duale de $J'$ se plonge
canoniquement dans $J$, donc hérite une
polarisation, d’où une polarisation}

On en conclut \inserted{par Néron-Tate} une forme quadratique
canonique \deleted{\unsure{sur}} définie positive \inserted{$Q$} sur $J'(k(\eta))$.

D’autre part, il y a \deleted{aussi} une forme
quadratique définie négative sur $\mathrm{NS}(X)'$
savoir la forme intersection.

\margin{$i\deleted{\unsure{\xi}}:\mathrm{NS}(X)'\longrightarrow J'(k(\eta))$}

À prouver la formule

\margin{?}

\[
 \boxed{Q(i(\unsure{\xi|\zeta}),i(\eta))=-\deleted{\unsure{\xi}}\xi\cdot\eta}
 \qquad(\xi,\eta\in\mathrm{NS}(X)')
\]

valable si les $Y_t$ sont tous \uline{intègres} ?
% archive sheet number: 14
